\documentclass[aspectratio=169]{beamer}

% Tema UFS --- Universidade Federal de Sergipe
\usetheme{UFS}

% Pacotes base
\usepackage[utf8]{inputenc}
\usepackage[portuguese]{babel}
\usepackage{amsmath,amssymb}
\usepackage{graphicx}
\usepackage{booktabs}
\usepackage{xcolor}

% TikZ e bibliotecas
\usepackage{tikz}
\usetikzlibrary{shapes.geometric, arrows.meta, positioning,
  matrix, fit, backgrounds, decorations.pathreplacing, calc}

% Listagens --- paleta VS Code Light+
\usepackage{listings}
\definecolor{bgcolor}{HTML}{FFFFFF}
\definecolor{linecolor}{HTML}{DDDDDD}
\definecolor{numcolor}{HTML}{999999}
\definecolor{kwcolor}{HTML}{0000FF}
\definecolor{strcolor}{HTML}{A31515}
\definecolor{cmtcolor}{HTML}{008000}

\lstdefinestyle{ufsStyle}{
  backgroundcolor=\color{bgcolor},
  basicstyle=\ttfamily\footnotesize\color{black},
  keywordstyle=\color{kwcolor}\bfseries,
  stringstyle=\color{strcolor},
  commentstyle=\color{cmtcolor}\itshape,
  numberstyle=\tiny\color{numcolor},
  numbers=left, numbersep=12pt,
  xleftmargin=20pt, framexleftmargin=20pt,
  breaklines=true, tabsize=4,
  keepspaces=true, showspaces=false,
  showstringspaces=false, showtabs=false,
  frame=single, rulecolor=\color{linecolor},
  framesep=4pt, captionpos=b, upquote=true,
  columns=fullflexible,
}
\lstset{style=ufsStyle, language=C}

\lstdefinestyle{tinyC}{
  style=ufsStyle, language=C, numbers=none,
  basicstyle=\ttfamily\scriptsize\color{black},
}
\lstdefinestyle{codC}{
  style=tinyC, basicstyle=\ttfamily\fontsize{7.5}{8.6}\selectfont\color{black},
}
\lstdefinestyle{codCs}{
  style=tinyC, basicstyle=\ttfamily\fontsize{7}{7.9}\selectfont\color{black},
}
\lstdefinestyle{shellStyle}{
  style=ufsStyle, language=bash, numbers=none,
  basicstyle=\ttfamily\scriptsize\color{black},
}
\lstdefinestyle{shellS}{
  style=shellStyle, basicstyle=\ttfamily\fontsize{7}{7.9}\selectfont\color{black},
}

% --- Estilos TikZ reutilizados ---
\tikzset{
  byte/.style={rectangle, draw=UFSBlue, fill=white, minimum width=0.36cm,
    minimum height=0.46cm, inner sep=0pt, font=\scriptsize\ttfamily},
  byteb/.style={byte, fill=UFSBlue!15},
  byter/.style={byte, draw=UFSRed, fill=UFSRed!12},
  bytev/.style={byte, draw=UFSGray!50, fill=black!4, text=UFSGray!70},
  caixa/.style={rectangle, draw=UFSBlue, fill=UFSBlue!10, rounded corners,
    align=center, font=\small, inner sep=5pt},
  caixar/.style={caixa, draw=UFSRed, fill=UFSRed!10},
  caixag/.style={caixa, draw=UFSGray, fill=black!4},
  item/.style={rectangle, draw=UFSBlue, fill=UFSBlue!10, rounded corners,
    minimum width=10.4cm, minimum height=0.5cm, font=\small},
  leg/.style={font=\scriptsize, text=UFSGray},
  seta/.style={-{Stealth[length=1.8mm]}, thick, UFSBlue},
  setar/.style={-{Stealth[length=1.8mm]}, thick, UFSRed},
  setag/.style={-{Stealth[length=1.8mm]}, thick, UFSGray},
}

% Uma fileira de caracteres, um por caixa: \fileira{x}{y}{estilo}{a,b,c}
\newcommand{\fileira}[4]{%
  \foreach \c [count=\k from 0] in {#4} {
    \node[#3] at (#1+0.36*\k, #2) {\c};
  }
}

% Uma fileira de bytes em hexadecimal, mais larga: \fileirax{x}{y}{estilo}{55,3f}
\newcommand{\fileirax}[4]{%
  \foreach \c [count=\k from 0] in {#4} {
    \node[#3, minimum width=0.5cm] at (#1+0.5*\k, #2) {\c};
  }
}

% Circunflexo em codigo no meio do texto
\newcommand{\cir}{\textasciicircum}

% Referencia da fonte no rodape do frame
\newcommand{\fonte}[1]{\vfill\hfill{\scriptsize\textcolor{UFSGray}{Fonte: #1}}}

% --- Metadados ---
\title{Arquivos binários}
\subtitle{\texttt{fread}, \texttt{fwrite} e \texttt{fseek}}
\author{Prof.\ André Yoshiaki Kashiwabara}
\institute{DCOMP -- Universidade Federal de Sergipe\\
\texttt{andreyoshiaki@dcomp.ufs.br}}
\date{}

\begin{document}

\begin{frame}[plain]
  \titlepage
\end{frame}

\begin{frame}{O que você vai aprender hoje}
  \begin{center}
  \begin{tikzpicture}
    \node[item] (t1) at (0, 0)    {Texto ou bytes: o que se perde ao gravar um número como texto};
    \node[item] (t2) at (0,-0.64) {A memória direto para o disco: \texttt{fwrite} e o modo \texttt{"wb"}};
    \node[item] (t3) at (0,-1.28) {O que há dentro de um registro: \texttt{sizeof}, preenchimento e ordem dos bytes};
    \node[item] (t4) at (0,-1.92) {De volta para a memória: \texttt{fread} e o que ele devolve};
    \node[item] (t5) at (0,-2.56) {Direto ao registro $k$: \texttt{fseek} e \texttt{ftell}};
    \node[item] (t6) at (0,-3.20) {Alterar sem reescrever: o modo \texttt{"r+b"}};
    \node[item] (t7) at (0,-3.84) {Exercícios};
    \draw[seta] (t1) -- (t2);
    \draw[seta] (t2) -- (t3);
    \draw[seta] (t3) -- (t4);
    \draw[seta] (t4) -- (t5);
    \draw[seta] (t5) -- (t6);
    \draw[seta] (t6) -- (t7);
  \end{tikzpicture}
  \end{center}
\end{frame}

% ==================================================================
\section{A pergunta de hoje}
% ==================================================================

\begin{frame}{O problema: achar e alterar um aluno no meio do arquivo}
  \vspace{-0.4em}
  \begin{columns}[T]
    \begin{column}{0.5\textwidth}
      \textbf{Até aqui}\\[0.3em]
      \small
      O \texttt{boletim\_csv.c} lê \texttt{turma.csv} linha a linha. Para
      chegar ao aluno 3\,000, lê as 2\,999 linhas anteriores. Para trocar
      uma nota, reescreve o arquivo inteiro --- \texttt{7.5} virar
      \texttt{10.0} muda o tamanho da linha.
    \end{column}
    \begin{column}{0.5\textwidth}
      \textbf{O que vamos descobrir}\\[0.3em]
      \small
      Se cada aluno ocupar \textbf{sempre o mesmo número de bytes}, o
      arquivo vira um vetor em disco: o aluno $k$ está numa posição que se
      calcula, e trocar uma nota é gravar por cima, sem mexer no resto.
    \end{column}
  \end{columns}
  \medskip
  \begin{block}{A pergunta de hoje}
    Como gravar uma \texttt{struct aluno} no disco exatamente como ela está
    na memória, ir direto ao aluno $k$ e alterá-lo ali mesmo? E o que se
    perde --- ou se ganha --- quando o arquivo deixa de ser texto?
  \end{block}
  \fonte{Deitel e Deitel (2011), cap.~11~\cite{deitel2011como};
    King (2008), §22.6--22.7~\cite{king2008c}}
\end{frame}

\begin{frame}[fragile]{Primeiro: o mesmo \texttt{double}, de dois jeitos}
  \vspace{-0.9em}
  \begin{columns}[T]
    \begin{column}{0.56\textwidth}
\begin{lstlisting}[style=codCs]
/* precisao.c (trecho) */
double x = 1.0 / 3.0, lido_t, lido_b;

FILE *ft = fopen("x.txt", "w");
FILE *fb = fopen("x.bin", "wb");
...  /* testa os dois NULL */
fprintf(ft, "%f", x);          /* texto */
fwrite(&x, sizeof x, 1, fb);   /* bytes */
fclose(ft);
fclose(fb);

...  /* reabre com "r" e "rb" */
fscanf(ft, "%lf", &lido_t);
fread(&lido_b, sizeof lido_b, 1, fb);

printf("texto:   %.17f  %s\n", lido_t,
       lido_t == x ? "igual" : "diferente");
printf("binario: %.17f  %s\n", lido_b,
       lido_b == x ? "igual" : "diferente");
\end{lstlisting}
    \end{column}
    \begin{column}{0.40\textwidth}
      \begin{block}{Antes de executar}
        \small
        \begin{enumerate}
          \item Quantos bytes terá \texttt{x.txt}? E \texttt{x.bin}?
          \item Os dois valores lidos de volta são iguais a \texttt{x}?
          \item O que \texttt{cat x.bin} mostra na tela?
        \end{enumerate}
      \end{block}
      \small \texttt{fwrite} grava os bytes que estão no endereço
      \texttt{\&x}; \texttt{fread} os copia de volta. Os detalhes vêm a
      seguir.
    \end{column}
  \end{columns}
\end{frame}

\begin{frame}[fragile]{Execute e compare}
\begin{lstlisting}[style=shellStyle]
$ gcc -Wall -Wextra precisao.c -o precisao && ./precisao
texto:   0.33333299999999999  diferente
binario: 0.33333333333333331  igual
$ wc -c x.txt x.bin
       8 x.txt
       8 x.bin
$ cat x.bin
UUUUUU?
\end{lstlisting}
  \begin{columns}[T]
    \begin{column}{0.48\textwidth}
      \small Os dois arquivos têm 8 bytes, mas só um guardou o número.
      \texttt{\%f} escreve seis casas: o texto arredondou, e
      \texttt{fscanf} não tem como recuperar o que foi jogado fora.
    \end{column}
    \begin{column}{0.48\textwidth}
      \small \texttt{cat} trata os bytes de \texttt{x.bin} como caracteres:
      \texttt{0x55} é a letra \texttt{U}; \texttt{0x3f}, o \texttt{?}. O arquivo não foi feito para ser
      lido por gente.
    \end{column}
  \end{columns}
\end{frame}

\begin{frame}{A resposta: o texto converte, o binário copia}
  \begin{center}
  \begin{tikzpicture}
    \node[leg, anchor=east] at (-0.4,0) {\texttt{x} na memória:};
    \fileirax{0}{0}{byteb}{55,55,55,55,55,55,d5,3f}
    \draw[seta] (1.75,-0.3) -- (1.75,-0.9) node[midway, right, leg] {\texttt{fwrite}: copia};
    \node[leg, anchor=east] at (-0.4,-1.2) {\texttt{x.bin}:};
    \fileirax{0}{-1.2}{byteb}{55,55,55,55,55,55,d5,3f}
    \draw[seta] (3.9,0) -- (5.7,0) node[midway, above, leg] {\texttt{fprintf "\%f"}};
    \node[leg, anchor=west] at (5.8,0.42) {\texttt{x.txt}: converte e arredonda};
    \fileira{6.0}{0}{byte}{0,.,3,3,3,3,3,3}
    \node[leg, text=UFSRed, anchor=west] at (8.8,0) {e o resto?};
  \end{tikzpicture}
  \end{center}
  \vspace{-0.3em}
  \begin{columns}[T]
    \begin{column}{0.5\textwidth}
      \begin{block}{Arquivo binário}
        \small Uma sequência de bytes sem linhas nem conversão: o que estava
        na memória vai para o disco tal e qual, e volta tal e qual.
      \end{block}
    \end{column}
    \begin{column}{0.46\textwidth}
      \small Dá para não perder nada em texto --- \texttt{\%.17g} escreve
      \texttt{0.33333333333333331}, 19 caracteres. Em binário,
      \textbf{todo} \texttt{double} tem 8 bytes e nenhuma conversão.
    \end{column}
  \end{columns}
  \fonte{ISO/IEC 9899:2018, §7.21.2~\cite{iso2018c17};
    Goldberg (1991), §2.1~\cite{goldberg1991what}}
\end{frame}

% ==================================================================
\section{Gravar bytes: \texttt{fwrite}}
% ==================================================================

\begin{frame}[fragile]{Definição: \texttt{fwrite} e \texttt{fread}}
  \vspace{-0.3em}
\begin{lstlisting}[style=codC]
size_t fwrite(const void *p, size_t tam, size_t n, FILE *f);
size_t fread(void *p, size_t tam, size_t n, FILE *f);
\end{lstlisting}
  \begin{columns}[T]
    \begin{column}{0.52\textwidth}
      \small \texttt{fwrite} grava \texttt{n} elementos de \texttt{tam} bytes
      cada, a partir do endereço \texttt{p}. \texttt{fread} lê \texttt{n}
      elementos e os coloca a partir de \texttt{p}. O \texttt{void *} aceita
      endereço de qualquer tipo: \texttt{\&x}, um vetor, uma
      \texttt{struct}.

      \medskip
      As duas devolvem quantos \textbf{elementos} --- não bytes ---
      conseguiram transferir.
    \end{column}
    \begin{column}{0.44\textwidth}
\begin{lstlisting}[style=codC]
double x;
int v[100];
fwrite(&x, sizeof x, 1, f);
fwrite(v, sizeof v[0], 100, f);
\end{lstlisting}
      \begin{block}{O \texttt{b} do modo}
        \small \texttt{"wb"}, \texttt{"rb"}: fluxo binário. No Linux e no
        macOS nada muda; no Windows, sem o \texttt{b}, cada
        \texttt{'\textbackslash n'} gravado vira dois bytes.
      \end{block}
    \end{column}
  \end{columns}
  \fonte{ISO/IEC 9899:2018, §7.21.2 e §7.21.8~\cite{iso2018c17};
    King (2008), §22.6~\cite{king2008c}}
\end{frame}

\begin{frame}[fragile]{O programa de hoje: \texttt{cadastro.c}}
  \vspace{-0.8em}
  \begin{columns}[T]
    \begin{column}{0.36\textwidth}
\begin{lstlisting}[style=codCs]
/* aluno.h */
#define NAV 4

struct aluno {
    int    mat;
    char   nome[50];
    double v[NAV];
};
\end{lstlisting}
      \small Converte \texttt{turma.csv}, da aula passada, em
      \texttt{turma.bin}: um registro \texttt{struct aluno} por aluno, um
      atrás do outro.

      \medskip
      A abertura é a de sempre, com \texttt{"r"} e \texttt{"wb"}.
    \end{column}
    \begin{column}{0.62\textwidth}
\begin{lstlisting}[style=codCs]
    char linha[256];
    int n = 0;
    while (fgets(linha, sizeof linha, ent) != NULL) {
        struct aluno a = {0};   /* comeca zerado */
        if (sscanf(linha, "%d;%49[^;];%lf;%lf;%lf;%lf",
                   &a.mat, a.nome, &a.v[0], &a.v[1],
                   &a.v[2], &a.v[3]) != 6)
            continue;
        if (fwrite(&a, sizeof a, 1, sai) != 1) {
            perror(nome_sai);
            break;
        }
        n++;
    }
    fclose(ent);
    fclose(sai);
    printf("%d alunos gravados em %s\n", n, nome_sai);
\end{lstlisting}
    \end{column}
  \end{columns}
  \fonte{King (2008), §22.6~\cite{king2008c}}
\end{frame}

\begin{frame}{O que este programa grava?}
  \begin{block}{Antes de compilar, escreva sua resposta}
    \begin{enumerate}
      \item Um \texttt{int} tem 4 bytes e um \texttt{double}, 8. Quantos
        bytes terá \texttt{turma.bin}, com quatro alunos?
      \item O que aparece com \texttt{strings turma.bin}, que mostra só os
        trechos legíveis?
      \item Onde, dentro do arquivo, estará o texto \texttt{Ana Souza}?
      \item Rodando duas vezes, o arquivo fica com quatro alunos ou com oito?
      \item Por que o \texttt{fwrite} recebe \texttt{\&a} e não \texttt{a}?
    \end{enumerate}
  \end{block}
\end{frame}

\begin{frame}[fragile]{Compile, execute e compare}
\begin{lstlisting}[style=shellStyle]
$ gcc -Wall -Wextra cadastro.c -o cadastro
$ ./cadastro
4 alunos gravados em turma.bin
$ wc -c turma.bin
     352 turma.bin
$ strings turma.bin       # so os trechos que parecem texto
Ana Souza
Bruno Lima
Carla Menezes Santos
Davi Rocha
\end{lstlisting}
  \begin{columns}[T]
    \begin{column}{0.48\textwidth}
      \small A conta $4 + 50 + 4 \times 8 = 86$ dá $4 \times 86 = 344$
      bytes. O arquivo tem \textbf{352}: sobram 8, dois por aluno.
    \end{column}
    \begin{column}{0.48\textwidth}
      \small Só os nomes aparecem: um \texttt{char} já é um caractere.
      Matrículas e notas são bytes que não formam texto. Rodar de novo deixa 352 bytes:
      \texttt{"wb"} apaga, como \texttt{"w"}.
    \end{column}
  \end{columns}
\end{frame}

\begin{frame}[fragile]{Onde estão os bytes a mais?}
  \vspace{-0.4em}
  \texttt{xxd} mostra cada byte em hexadecimal; à direita, os que são letras.
\begin{lstlisting}[style=shellStyle, language={}]
$ xxd -l 88 turma.bin           # so os primeiros 88 bytes
00000000: a56e 130c 416e 6120 536f 757a 6100 0000  .n..Ana Souza...
00000010: 0000 0000 0000 0000 0000 0000 0000 0000  ................
00000020: 0000 0000 0000 0000 0000 0000 0000 0000  ................
00000030: 0000 0000 0000 0000 0000 0000 0000 2040  .............. @
00000040: 0000 0000 0000 1a40 0000 0000 0000 2240  .......@......"@
00000050: 0000 0000 0000 1e40                      .......@
\end{lstlisting}
  \begin{block}{Investigue}
    \small
    \begin{enumerate}
      \item Os primeiros 88 bytes são o primeiro aluno. Em que posição (em
        hexadecimal) começa o nome? E onde ele termina?
      \item O \texttt{double} \texttt{8.0} termina com os bytes
        \texttt{20 40}. Em que posição começa \texttt{v[0]}? Por que não
        logo depois do nome?
    \end{enumerate}
  \end{block}
\end{frame}

\begin{frame}{A resposta: \texttt{sizeof(struct aluno)} é 88}
  \begin{center}
  \begin{tikzpicture}[x=0.125cm]
    \draw[fill=UFSBlue!25, draw=UFSBlue] (0,0) rectangle (4,0.6);
    \node[font=\tiny\ttfamily] at (2,0.3) {mat};
    \draw[fill=UFSBlue!8, draw=UFSBlue] (4,0) rectangle (54,0.6);
    \node[font=\scriptsize\ttfamily] at (29,0.3) {nome[50]};
    \draw[fill=UFSRed!25, draw=UFSRed] (54,0) rectangle (56,0.6);
    \foreach \i in {0,...,3} {
      \draw[fill=UFSBlue!25, draw=UFSBlue] (56+8*\i,0) rectangle (64+8*\i,0.6);
      \node[font=\tiny\ttfamily] at (60+8*\i,0.3) {v[\i]};
    }
    \foreach \p in {0,4,64,72,80,88}
      \node[leg, font=\tiny] at (\p,-0.25) {\p};
    \node[leg, font=\tiny, anchor=east] at (54.3,-0.25) {54};
    \node[leg, font=\tiny, anchor=west] at (55.7,-0.25) {56};
    \draw[setar] (55,1.25) -- (55,0.7);
    \node[leg, text=UFSRed, anchor=south] at (55,1.2) {2 bytes de preenchimento};
  \end{tikzpicture}
  \end{center}
  \vspace{-0.4em}
  \begin{columns}[T]
    \begin{column}{0.5\textwidth}
      \small O processador lê um \texttt{double} mais depressa --- em alguns,
      só consegue lê-lo --- se o endereço for múltiplo de 8. O compilador
      então deixa \texttt{nome} terminar em 54 e começa \texttt{v} em 56,
      com dois bytes sem nome no meio.
    \end{column}
    \begin{column}{0.46\textwidth}
      \begin{block}{Regra}
        \small O tamanho de um registro é \texttt{sizeof(struct aluno)},
        nunca a soma dos campos feita à mão. O compilador pode pôr
        preenchimento entre os campos e no fim --- só não no início.
      \end{block}
    \end{column}
  \end{columns}
  \fonte{ISO/IEC 9899:2018, §6.7.2.1~\cite{iso2018c17};
    Bryant e O'Hallaron (2016), §3.9.3~\cite{bryant2016computer}}
\end{frame}

\begin{frame}{E a matrícula? Os bytes estão de trás para frente}
  \begin{center}
  \begin{tikzpicture}
    \node[leg, anchor=east] at (-0.4,0) {no arquivo, posições 0--3:};
    \foreach \b [count=\k from 0] in {a5,6e,13,0c}
      \node[byteb, minimum width=0.7cm] (m\k) at (0.7*\k,0) {\b};
    \node[leg, anchor=east] at (-0.4,-1.1) {o número \texttt{202600101}:};
    \foreach \b [count=\k from 0] in {0c,13,6e,a5}
      \node[byte, minimum width=0.7cm] (n\k) at (0.7*\k,-1.1) {\b};
    \node[leg, anchor=west] at (2.6,-1.1) {$= \mathtt{0x0c136ea5}$};
    \draw[setag] (m0.south) -- (n3.north);
    \draw[setag] (m3.south) -- (n0.north);
    \node[leg, anchor=west] at (2.6,0) {byte menos significativo primeiro};
  \end{tikzpicture}
  \end{center}
  \vspace{-0.3em}
  \begin{columns}[T]
    \begin{column}{0.48\textwidth}
      \small Os PCs (x86) e os Macs com processador Apple guardam o byte
      \textbf{menos} significativo no endereço menor: \emph{little endian}.
      Outras máquinas fazem o contrário, \emph{big endian}. O
      \texttt{fwrite} grava a ordem da máquina em que roda.
    \end{column}
    \begin{column}{0.48\textwidth}
      \begin{alertblock}{O arquivo binário depende da máquina}
        \small Ordem dos bytes, tamanho de \texttt{int} e \texttt{long},
        preenchimento: tudo isso vai para o disco. Ler o \texttt{turma.bin}
        em outra arquitetura ou com outro compilador pode dar lixo.
      \end{alertblock}
    \end{column}
  \end{columns}
  \fonte{Bryant e O'Hallaron (2016), §2.1.3~\cite{bryant2016computer};
    ISO/IEC 9899:2018, §6.2.6.1~\cite{iso2018c17}}
\end{frame}

\begin{frame}{Recapitulando: \texttt{fwrite}}
  \begin{block}{O que aprendemos}
    \begin{itemize}
      \item \texttt{fwrite(p, tam, n, f)} copia \texttt{n $\times$ tam}
        bytes da memória para o arquivo, sem conversão;
      \item o modo leva \texttt{b}: \texttt{"wb"} apaga e grava, como
        \texttt{"w"};
      \item devolve quantos elementos gravou: compare com \texttt{n};
      \item o tamanho do registro é \texttt{sizeof}, que inclui o
        preenchimento;
      \item o arquivo guarda a ordem dos bytes e o \emph{layout} da máquina
        que o gravou.
    \end{itemize}
  \end{block}
\end{frame}

% ==================================================================
\section{De volta: \texttt{fread}}
% ==================================================================

\begin{frame}[fragile]{\texttt{lista.c}: um registro por volta}
  \vspace{-0.9em}
  \begin{columns}[T]
    \begin{column}{0.56\textwidth}
\begin{lstlisting}[style=codCs]
/* lista.c (trecho) */
int main(int argc, char *argv[])
{
    const char *nome = argc > 1 ? argv[1]
                                : "turma.bin";
    FILE *f = fopen(nome, "rb");
    if (f == NULL) {
        perror(nome);
        return 1;
    }
    struct aluno a;
    while (fread(&a, sizeof a, 1, f) == 1) {
        double m = (a.v[0] + a.v[1] + a.v[2]
                    + a.v[3]) / NAV;
        printf("%d %-22s %5.2f\n", a.mat, a.nome, m);
    }
    fclose(f);
    return 0;
}
\end{lstlisting}
    \end{column}
    \begin{column}{0.40\textwidth}
\begin{lstlisting}[style=shellS]
$ ./lista
202600101 Ana Souza    7.75
202600102 Bruno Lima   5.75
202600103 Carla Me...  9.25
202600104 Davi Rocha   4.00
\end{lstlisting}
      \small Nenhum \texttt{sscanf}, nenhum separador: cada
      \texttt{fread} enche os 88 bytes de \texttt{a} e os campos já estão
      prontos.

      \medskip
      O laço para quando \texttt{fread} não consegue um registro
      inteiro e devolve \texttt{0} --- o mesmo cuidado do
      \texttt{fscanf(...) == 5}.
    \end{column}
  \end{columns}
  \fonte{King (2008), §22.6~\cite{king2008c}}
\end{frame}

\begin{frame}[fragile]{Vários de uma vez: a turma inteira num vetor}
  \vspace{-0.4em}
  \begin{columns}[T]
    \begin{column}{0.58\textwidth}
\begin{lstlisting}[style=codCs]
/* n registros cabem num malloc so */
struct aluno *t = malloc(n * sizeof *t);
if (t == NULL) { ... }
size_t lidos = fread(t, sizeof *t, n, f);
if (lidos < n)
    fprintf(stderr, "so %zu de %zu\n", lidos, n);
...
fwrite(t, sizeof *t, lidos, g);  /* tudo de volta */
free(t);
\end{lstlisting}
      \small Um \texttt{fread} com \texttt{n} elementos substitui um laço de
      \texttt{n} leituras. No fim do arquivo, ele devolve quantos
      \emph{conseguiu}, que pode ser menos que \texttt{n}.
    \end{column}
    \begin{column}{0.38\textwidth}
      \begin{block}{Bytes como elementos}
        \small Com \texttt{tam = 1}, o valor devolvido é o número de
        bytes:
      \end{block}
\begin{lstlisting}[style=codCs]
char buf[4096];
size_t k;
while ((k = fread(buf, 1,
           sizeof buf, f)) > 0)
    fwrite(buf, 1, k, g);
\end{lstlisting}
      \small É assim que se copia qualquer arquivo --- imagem, música,
      executável.
    \end{column}
  \end{columns}
  \fonte{ISO/IEC 9899:2018, §7.21.8.1~\cite{iso2018c17};
    Kernighan e Ritchie (1988), §B1.5~\cite{kernighan1988c}}
\end{frame}

\begin{frame}[fragile]{E se a \texttt{struct} tiver um ponteiro?}
  \vspace{-0.9em}
  \begin{columns}[T]
    \begin{column}{0.56\textwidth}
\begin{lstlisting}[style=codCs]
/* ponteiro.c (trecho) */
struct contato {
    int   mat;
    char *nome;        /* ponteiro, nao vetor */
};
...
if (argc > 1 && strcmp(argv[1], "grava") == 0) {
    c.mat = 202600101;
    c.nome = malloc(50);
    strcpy(c.nome, "Ana Souza");
    ...  /* abre contato.bin com "wb" */
    fwrite(&c, sizeof c, 1, f);
    printf("gravei %d %s (%zu bytes)\n",
           c.mat, c.nome, sizeof c);
} else {
    ...  /* abre contato.bin com "rb" */
    fread(&c, sizeof c, 1, f);
    printf("li %d\n", c.mat);
    printf("nome: %s\n", c.nome);
}
\end{lstlisting}
    \end{column}
    \begin{column}{0.40\textwidth}
\begin{lstlisting}[style=shellS]
$ ./ponteiro grava
$ ./ponteiro le
\end{lstlisting}
      \begin{block}{Antes de executar}
        \small
        \begin{enumerate}
          \item Quantos bytes \texttt{sizeof c} informa? O nome cabe neles?
          \item A segunda execução imprime a matrícula? E o nome?
        \end{enumerate}
      \end{block}
    \end{column}
  \end{columns}
\end{frame}

\begin{frame}[fragile]{Execute e compare}
\begin{lstlisting}[style=shellStyle]
$ ./ponteiro grava
gravei 202600101 Ana Souza (16 bytes)
$ ./ponteiro le
li 202600101
Segmentation fault (core dumped)
$ xxd contato.bin
00000000: a56e 130c ffff ffff 909b 8d04 0100 0000  .n..............
\end{lstlisting}
  \begin{columns}[T]
    \begin{column}{0.48\textwidth}
      \begin{block}{O nome não está no arquivo}
        \small Os 16 bytes são a matrícula, 4 bytes de preenchimento --- com
        o lixo que havia na pilha --- e os 8 bytes de um
        \textbf{endereço}.
      \end{block}
    \end{column}
    \begin{column}{0.48\textwidth}
      \small Na segunda execução, aquele endereço não pertence a mais
      ninguém: o \texttt{malloc} que o criou morreu com o primeiro processo.
      O \texttt{printf} seguiu o ponteiro e caiu.
    \end{column}
  \end{columns}
\end{frame}

\begin{frame}{A resposta: \texttt{fwrite} copia o ponteiro, não o que ele aponta}
  \begin{center}
  \begin{tikzpicture}
    \node[leg] at (1.2,0.75) {\texttt{c}, na memória};
    \node[byteb, minimum width=1.2cm] (cm) at (0.6,0) {mat};
    \node[byteb, minimum width=1.2cm] (cn) at (1.8,0) {nome};
    \fileira{4.2}{0}{byte}{A,n,a}
    \node[byte] at (5.28,0) {\textvisiblespace};
    \fileira{5.64}{0}{byte}{S,o,u,z,a}
    \node[byte, minimum width=0.5cm] at (7.6,0) {\textbackslash 0};
    \node[leg] at (5.9,0.75) {bloco do \texttt{malloc}};
    \draw[seta] (cn.east) -- (4.0,0);
    \draw[decorate, decoration={brace, amplitude=4pt, mirror}, thick, UFSRed]
      (0,-0.35) -- (2.4,-0.35)
      node[midway, below=5pt, leg, text=UFSRed] {só isto vai para o disco};
  \end{tikzpicture}
  \end{center}
  \vspace{-0.3em}
  \begin{columns}[T]
    \begin{column}{0.48\textwidth}
      \small \texttt{fwrite} não sabe o que os bytes significam: copia
      \texttt{sizeof c} bytes a partir de \texttt{\&c}, e o campo
      \texttt{nome} é só um número de 8 bytes. Um endereço só vale dentro do
      processo que o recebeu.
    \end{column}
    \begin{column}{0.48\textwidth}
      \begin{block}{Para gravar com \texttt{fwrite}}
        \small A \texttt{struct} precisa ser ``plana'': vetores de tamanho
        fixo, como \texttt{char nome[50]}, e nenhum ponteiro. Uma lista
        encadeada se grava nó a nó, só com os dados --- os
        \texttt{prox} se refazem ao ler.
      \end{block}
    \end{column}
  \end{columns}
  \fonte{Kernighan e Ritchie (1988), §6.4~\cite{kernighan1988c};
    King (2008), §22.6~\cite{king2008c}}
\end{frame}

\begin{frame}{Recapitulando: \texttt{fread}}
  \begin{block}{O que aprendemos}
    \begin{itemize}
      \item \texttt{fread(p, tam, n, f)} enche \texttt{n} elementos a
        partir de \texttt{p} e devolve quantos conseguiu;
      \item no laço, teste \texttt{== 1}: no fim, \texttt{fread} devolve 0;
      \item um vetor inteiro vai e volta com uma chamada só;
      \item com \texttt{tam = 1}, o retorno conta bytes --- serve para
        copiar qualquer arquivo;
      \item ponteiros não sobrevivem ao disco: grave vetores, não
        \texttt{char *}.
    \end{itemize}
  \end{block}
\end{frame}

% ==================================================================
\section{Direto ao registro: \texttt{fseek}}
% ==================================================================

\begin{frame}{Registros do mesmo tamanho: o arquivo é um vetor}
  \begin{center}
  \begin{tikzpicture}[x=0.125cm]
    \foreach \k/\n in {0/Ana, 1/Bruno, 2/Carla, 3/Davi} {
      \draw[fill=UFSBlue!10, draw=UFSBlue] (22*\k,0) rectangle (22*\k+22,0.7);
      \node[font=\scriptsize] at (22*\k+11,0.35) {registro \k: \n};
      \node[leg, font=\tiny] at (22*\k,-0.25) {\pgfmathparse{int(88*\k)}\pgfmathresult};
    }
    \node[leg, font=\tiny] at (88,-0.25) {352};
    \draw[setar] (44,1.35) -- (44,0.8);
    \node[leg, text=UFSRed, anchor=south] at (44,1.3) {registro 2 começa em $2 \times 88 = 176$};
    \node[leg, anchor=west] at (90,0.35) {fim};
  \end{tikzpicture}
  \end{center}
  \vspace{-0.3em}
  \begin{columns}[T]
    \begin{column}{0.48\textwidth}
      \small Em \texttt{turma.bin}, o registro $k$ começa no byte
      $k \times \mathtt{sizeof(struct\ aluno)}$, contando de zero, como num
      vetor. O número de registros é o tamanho do arquivo dividido por 88.
    \end{column}
    \begin{column}{0.48\textwidth}
      \small Em \texttt{turma.csv} não há conta possível: cada linha tem um
      tamanho, e só se sabe onde começa a quarta depois de ler as três
      primeiras.
    \end{column}
  \end{columns}
  \fonte{Deitel e Deitel (2011), cap.~11~\cite{deitel2011como}}
\end{frame}

\begin{frame}[fragile]{Definição: \texttt{fseek}, \texttt{ftell} e \texttt{rewind}}
  \vspace{-0.3em}
\begin{lstlisting}[style=codC]
int  fseek(FILE *f, long desloc, int origem);   /* 0 se deu certo */
long ftell(FILE *f);                            /* posicao atual, em bytes */
void rewind(FILE *f);                           /* volta ao byte 0 */
\end{lstlisting}
  \begin{columns}[T]
    \begin{column}{0.5\textwidth}
      \begin{center}
      \small
      \begin{tabular}{@{}ll@{}}
        \toprule
        \texttt{origem} & \texttt{desloc} conta a partir de\\
        \midrule
        \texttt{SEEK\_SET} & o início do arquivo\\
        \texttt{SEEK\_CUR} & a posição atual\\
        \texttt{SEEK\_END} & o fim do arquivo\\
        \bottomrule
      \end{tabular}
      \end{center}
      \small \texttt{fseek} só move a posição que o \texttt{FILE} guarda; a
      próxima leitura ou gravação começa ali.
    \end{column}
    \begin{column}{0.46\textwidth}
\begin{lstlisting}[style=codCs]
fseek(f, k * (long) sizeof a, SEEK_SET);
fseek(f, -(long) sizeof a, SEEK_CUR);
fseek(f, 0, SEEK_END);
long tam = ftell(f);
\end{lstlisting}
      \small O deslocamento é \texttt{long} e pode ser negativo:
      \texttt{sizeof} é \texttt{size\_t}, sem sinal --- converta
      \emph{antes} de pôr o sinal de menos.
    \end{column}
  \end{columns}
  \fonte{ISO/IEC 9899:2018, §7.21.9~\cite{iso2018c17};
    King (2008), §22.7~\cite{king2008c}}
\end{frame}

\begin{frame}[fragile]{\texttt{consulta.c}: vai direto ao registro $k$}
  \vspace{-0.6em}
  \begin{columns}[T]
    \begin{column}{0.6\textwidth}
\begin{lstlisting}[style=codCs]
/* consulta.c (trecho): ./consulta turma.bin k */
fseek(f, 0, SEEK_END);              /* vai ao fim... */
long n = ftell(f) / (long) sizeof(struct aluno);

long k = atol(argv[2]);
struct aluno a;
fseek(f, k * (long) sizeof a, SEEK_SET);
if (fread(&a, sizeof a, 1, f) != 1) {
    printf("registro %ld nao existe "
           "(o arquivo tem %ld)\n", k, n);
    fclose(f);
    return 1;
}
printf("registro %ld de %ld: %d %s\n",
       k, n, a.mat, a.nome);
\end{lstlisting}
    \end{column}
    \begin{column}{0.36\textwidth}
      \begin{block}{Antes de executar}
        \small
        \begin{enumerate}
          \item Quem aparece com \texttt{./consulta turma.bin 1}?
          \item E com \texttt{4}? O \texttt{fseek} para além do fim
            falha?
          \item E com \texttt{-1}?
          \item Quantos \texttt{fread} o programa faz para chegar ao
            registro 3?
        \end{enumerate}
      \end{block}
    \end{column}
  \end{columns}
\end{frame}

\begin{frame}[fragile]{Execute e compare}
  \begin{columns}[T]
    \begin{column}{0.54\textwidth}
\begin{lstlisting}[style=shellS]
$ ./consulta turma.bin 1
registro 1 de 4: 202600102 Bruno Lima
$ ./consulta turma.bin 3
registro 3 de 4: 202600104 Davi Rocha
$ ./consulta turma.bin 4
registro 4 nao existe (o arquivo tem 4)
$ ./consulta turma.bin -1
registro -1 nao existe (o arquivo tem 4)
\end{lstlisting}
    \end{column}
    \begin{column}{0.42\textwidth}
      \small O registro 1 é o \textbf{segundo} aluno: conta-se de zero. E
      chegar ao registro 3 custou um \texttt{fread} só --- com três mil
      alunos também.
    \end{column}
  \end{columns}
  \medskip
  \begin{columns}[T]
    \begin{column}{0.48\textwidth}
      \begin{block}{Com \texttt{k = 4}}
        \small O \texttt{fseek} para o byte 352 \textbf{deu certo}: posicionar
        além do fim é permitido. Quem falhou foi o \texttt{fread}, que não
        encontrou bytes ali.
      \end{block}
    \end{column}
    \begin{column}{0.48\textwidth}
      \begin{block}{Com \texttt{k = -1}}
        \small Antes do início, o \texttt{fseek} falha e devolve $-1$; a
        posição continua no fim, onde o primeiro \texttt{fseek} a deixou, e
        o \texttt{fread} também falha.
      \end{block}
    \end{column}
  \end{columns}
\end{frame}

\begin{frame}{A resposta: \texttt{fseek} posiciona, quem lê decide}
  \begin{center}
  \begin{tikzpicture}[x=0.11cm]
    \foreach \k in {0,...,3}
      \draw[fill=UFSBlue!10, draw=UFSBlue] (22*\k,0) rectangle (22*\k+22,0.6);
    \draw[dashed, draw=UFSGray] (88,0) rectangle (110,0.6);
    \node[leg] at (99,0.3) {nada};
    \foreach \k/\p in {1/22, 4/88} {
      \draw[seta] (\p,1.15) -- (\p,0.7);
      \node[leg, anchor=south] at (\p,1.1) {\texttt{k = \k}};
    }
    \draw[setar] (-6,1.15) -- (-6,0.7);
    \node[leg, text=UFSRed, anchor=south] at (-6,1.1) {\texttt{k = -1}: recusado};
  \end{tikzpicture}
  \end{center}
  \vspace{-0.3em}
  \begin{columns}[T]
    \begin{column}{0.48\textwidth}
      \small Teste sempre o que o \texttt{fread} devolve depois de um
      \texttt{fseek} --- é ele que diz se havia um registro ali. Se for
      \emph{gravar} depois de saltar além do fim, o arquivo cresce, e o
      intervalo pulado vira bytes zero.
    \end{column}
    \begin{column}{0.48\textwidth}
      \begin{block}{Ler sequencial ou saltar?}
        \small \texttt{fseek} + \texttt{fread}: uma leitura, seja qual for
        $k$. Ler do início até $k$: $k+1$ leituras. É a diferença entre
        \texttt{t[k]} e percorrer uma lista encadeada.
      \end{block}
    \end{column}
  \end{columns}
  \fonte{ISO/IEC 9899:2018, §7.21.9.2~\cite{iso2018c17};
    Kerrisk (2010), §4.7~\cite{kerrisk2010linux}}
\end{frame}

\begin{frame}{Recapitulando: \texttt{fseek} e \texttt{ftell}}
  \begin{block}{O que aprendemos}
    \begin{itemize}
      \item com registros de tamanho fixo, o registro $k$ começa em
        $k \times \mathtt{sizeof}$;
      \item \texttt{fseek(f, desloc, origem)} move a posição;
        \texttt{ftell(f)} a informa;
      \item \texttt{fseek(f, 0, SEEK\_END)} + \texttt{ftell} dão o tamanho
        do arquivo;
      \item posicionar além do fim é permitido: confira o retorno do
        \texttt{fread};
      \item converta \texttt{sizeof} para \texttt{long} antes de
        multiplicar ou negar.
    \end{itemize}
  \end{block}
\end{frame}

% ==================================================================
\section{Alterar sem reescrever}
% ==================================================================

\begin{frame}[fragile]{\texttt{corrige.c}: trocar uma nota no lugar}
  \vspace{-0.6em}
  \begin{columns}[T]
    \begin{column}{0.55\textwidth}
\begin{lstlisting}[style=codCs]
/* corrige.c (trecho) */
FILE *f = fopen(argv[1], "r+b");  /* ler e gravar */
if (f == NULL) { ... }
struct aluno a;
while (fread(&a, sizeof a, 1, f) == 1) {
    if (a.mat != mat)
        continue;
    a.v[av] = nota;
    fseek(f, -(long) sizeof a, SEEK_CUR);
    if (fwrite(&a, sizeof a, 1, f) != 1)
        perror(argv[1]);
    else
        printf("%d: nota %d agora vale %.1f\n",
               mat, av, nota);
    fclose(f);
    return 0;
}
printf("matricula %d nao encontrada\n", mat);
\end{lstlisting}
    \end{column}
    \begin{column}{0.43\textwidth}
      \small \texttt{"r+b"} abre para ler \textbf{e} gravar, sem apagar o
      conteúdo.

      \medskip
      Depois do \texttt{fread}, a posição está no fim do registro lido. O
      \texttt{fseek} recua 88 bytes, para o \texttt{fwrite} cair em cima do
      mesmo aluno.
\begin{lstlisting}[style=shellS]
$ ./corrige turma.bin 202600104 3 6.0
202600104: nota 3 agora vale 6.0
$ ./lista | tail -1
202600104 Davi Rocha   4.88
$ wc -c turma.bin
     352 turma.bin
\end{lstlisting}
    \end{column}
  \end{columns}
  \fonte{ISO/IEC 9899:2018, §7.21.5.3~\cite{iso2018c17};
    Deitel e Deitel (2011), cap.~11~\cite{deitel2011como}}
\end{frame}

\begin{frame}[fragile]{E se o arquivo for aberto com \texttt{"rb"}?}
  \vspace{-0.6em}
  \begin{columns}[T]
    \begin{column}{0.56\textwidth}
\begin{lstlisting}[style=codCs]
/* corrige_rb.c: "rb" e fwrite sem teste */
FILE *f = fopen(argv[1], "rb");
if (f == NULL) { ... }
struct aluno a;
while (fread(&a, sizeof a, 1, f) == 1) {
    if (a.mat != mat)
        continue;
    a.v[av] = nota;
    fseek(f, -(long) sizeof a, SEEK_CUR);
    fwrite(&a, sizeof a, 1, f);
    printf("%d: nota %d agora vale %.1f\n",
           mat, av, nota);
    break;
}
\end{lstlisting}
    \end{column}
    \begin{column}{0.40\textwidth}
\begin{lstlisting}[style=shellS]
$ ./cadastro
$ ./corrige_rb turma.bin 202600104 3 6.0
$ ./lista | tail -1
\end{lstlisting}
      \begin{block}{Antes de executar}
        \small
        \begin{enumerate}
          \item O compilador reclama do \texttt{fwrite} com \texttt{"rb"}?
          \item A mensagem \texttt{nota 3 agora vale 6.0} aparece?
          \item Que média a \texttt{lista} mostra para Davi: 4.00 ou 4.88?
        \end{enumerate}
      \end{block}
    \end{column}
  \end{columns}
\end{frame}

\begin{frame}[fragile]{Execute e compare}
\begin{lstlisting}[style=shellStyle]
$ ./cadastro
4 alunos gravados em turma.bin
$ ./corrige_rb turma.bin 202600104 3 6.0
202600104: nota 3 agora vale 6.0
$ ./lista | tail -1
202600104 Davi Rocha              4.00
\end{lstlisting}
  \begin{columns}[T]
    \begin{column}{0.48\textwidth}
      \begin{block}{O programa mentiu}
        \small Compilou sem aviso, rodou sem erro e anunciou a correção. A
        nota não mudou: o \texttt{fwrite} falhou e devolveu 0, e ninguém
        olhou.
      \end{block}
    \end{column}
    \begin{column}{0.48\textwidth}
      \small O modo é verificado em tempo de execução, pelo \texttt{FILE},
      não pelo compilador. Com o retorno testado, \texttt{perror} teria
      dito \emph{Bad file descriptor}.
    \end{column}
  \end{columns}
\end{frame}

\begin{frame}{A resposta: os modos binários e a regra do \texttt{fseek}}
  \vspace{-0.3em}
  \begin{columns}[T]
    \begin{column}{0.52\textwidth}
      \begin{center}
      \scriptsize
      \begin{tabular}{@{}lll@{}}
        \toprule
        modo & faz & se existe\\
        \midrule
        \texttt{"rb"}  & só lê            & mantém\\
        \texttt{"wb"}  & só grava         & \textcolor{UFSRed}{apaga}\\
        \texttt{"ab"}  & grava no fim     & mantém\\
        \texttt{"r+b"} & lê e grava       & mantém; \texttt{NULL} se não existe\\
        \texttt{"w+b"} & lê e grava       & \textcolor{UFSRed}{apaga}\\
        \texttt{"a+b"} & lê; grava no fim & mantém\\
        \bottomrule
      \end{tabular}
      \end{center}
      \small Para alterar registros de um arquivo que já existe:
      \texttt{"r+b"}.
    \end{column}
    \begin{column}{0.44\textwidth}
      \begin{block}{Ao trocar de leitura para gravação}
        \small Num fluxo aberto com \texttt{+}, entre um \texttt{fread} e um
        \texttt{fwrite} é obrigatório um \texttt{fseek} (ou
        \texttt{rewind}); entre um \texttt{fwrite} e um \texttt{fread}, um
        \texttt{fseek} ou \texttt{fflush}. Sem isso, o comportamento é
        indefinido --- por causa do buffer.
      \end{block}
      \small E teste o retorno de \texttt{fwrite} como o de
      \texttt{fopen}.
    \end{column}
  \end{columns}
  \fonte{ISO/IEC 9899:2018, §7.21.5.3~\cite{iso2018c17};
    King (2008), §22.2~\cite{king2008c}}
\end{frame}

\begin{frame}{Recapitulando: texto ou binário?}
  \begin{center}
  \footnotesize
  \begin{tabular}{@{}lll@{}}
    \toprule
    & arquivo texto (\texttt{turma.csv}) & arquivo binário (\texttt{turma.bin})\\
    \midrule
    abre num editor          & sim                         & não: precisa de \texttt{xxd}\\
    um \texttt{double}       & caracteres, com arredondamento & 8 bytes, exato\\
    ler um aluno             & \texttt{fgets} + \texttt{sscanf} & um \texttt{fread}\\
    ir ao aluno $k$          & ler os $k$ anteriores       & \texttt{fseek} + \texttt{fread}\\
    trocar uma nota          & reescrever o arquivo        & \texttt{fseek} + \texttt{fwrite}\\
    outra máquina            & funciona                    & pode dar lixo\\
    \bottomrule
  \end{tabular}
  \end{center}
  \begin{block}{Qual escolher}
    \footnotesize Texto para trocar dados com pessoas e outros programas;
    binário para dados que só o seu programa lê, com acesso direto e
    registros de tamanho fixo.
  \end{block}
  \fonte{King (2008), §22.1~\cite{king2008c}}
\end{frame}

% ==================================================================
\section{Mãos à obra}
% ==================================================================

\begin{frame}{Altere o programa}
  \begin{block}{Tarefa 1 --- acrescentar alunos}
    \small Troque o modo do \texttt{cadastro.c} para \texttt{"ab"} e rode-o
    duas vezes. Quantos alunos a \texttt{lista} mostra? E o que o
    \texttt{consulta} responde para \texttt{k = 5}?
  \end{block}
  \begin{block}{Tarefa 2 --- onde está a posição}
    \small Na \texttt{lista}, imprima \texttt{ftell(f)} antes e depois de
    cada \texttt{fread}. Que números aparecem, e por que o primeiro é 0?
  \end{block}
  \begin{block}{Tarefa 3 --- um cabeçalho}
    \small Faça o \texttt{cadastro} gravar, antes dos alunos, um
    \texttt{int} com a quantidade --- grave 0, e no fim volte com
    \texttt{rewind} e grave o valor certo. Ajuste a \texttt{lista} e o
    \texttt{consulta} para pular esses 4 bytes.
  \end{block}
\end{frame}

\begin{frame}{Exercícios (1/3)}
  \footnotesize
  \begin{block}{Exercício 1 --- quadrados em disco}
    \texttt{./quadrados n saida.bin} grava os \texttt{int} $1, 4, 9,
    \ldots, n^2$ com \textbf{um único} \texttt{fwrite}: monte o vetor com
    \texttt{malloc} antes. Quantos bytes terá o arquivo para
    $n = 1\,000$? E se fosse texto, um número por linha?
  \end{block}
  \begin{block}{Exercício 2 --- copiar qualquer arquivo}
    \texttt{./copia origem destino} copia em blocos de 4\,096 bytes com
    \texttt{fread} e \texttt{fwrite} e imprime \texttt{10000 bytes
    copiados}. Teste com uma imagem e confira com \texttt{cmp}: o último
    bloco é menor que os outros.
  \end{block}
  \fonte{King (2008), §22.6~\cite{king2008c};
    Kernighan e Ritchie (1988), §B1.5~\cite{kernighan1988c}}
\end{frame}

\begin{frame}{Exercícios (2/3)}
  \footnotesize
  \begin{block}{Exercício 3 --- comparar dois arquivos}
    \texttt{./compara a.bin b.bin} imprime \texttt{iguais}, ou
    \texttt{diferem no byte 41}, ou \texttt{a.bin termina antes, no byte
    100} quando um arquivo é o começo do outro. Bytes contados de zero.
  \end{block}
  \begin{block}{Exercício 4 --- contar registros sem ler}
    \texttt{./registros turma.bin} imprime \texttt{4 registros de 88
    bytes} usando só \texttt{fseek} e \texttt{ftell}. Se o tamanho não for
    múltiplo de \texttt{sizeof(struct aluno)}, imprima
    \texttt{arquivo corrompido: 10 bytes sobrando}.
  \end{block}
  \fonte{King (2008), §22.6--22.7~\cite{king2008c}}
\end{frame}

\begin{frame}[fragile]{Exercícios (3/3)}
  \footnotesize
  \begin{columns}[T]
    \begin{column}{0.48\textwidth}
      \begin{block}{Exercício 5 --- de trás para frente}
        \texttt{./inverte turma.bin} lista os alunos do último ao primeiro,
        com \texttt{fseek}, sem guardar o arquivo na memória.
      \end{block}
      \begin{block}{Exercício 6 --- um \texttt{xxd} simples}
        \texttt{./hexa arquivo} mostra 16 bytes por linha: a posição, os
        bytes em hexadecimal e, entre barras, os que são imprimíveis
        (\texttt{isprint}); os outros viram um ponto.
      \end{block}
    \end{column}
    \begin{column}{0.48\textwidth}
\begin{lstlisting}[style=shellS, language={}]
$ ./inverte turma.bin
202600104 Davi Rocha
202600103 Carla Menezes Santos
202600102 Bruno Lima
202600101 Ana Souza
$ ./hexa ola.bin
00000000  4f 6c 61 2c 20 6d 75 6e
          64 6f 21 0a 00 01 ff 7f
          |Ola, mundo!.....|
\end{lstlisting}
      \scriptsize\raggedright No \texttt{hexa}, os 16 bytes e as barras
      ficam numa linha só; a quebra acima é só para caber no slide.
    \end{column}
  \end{columns}
  \fonte{King (2008), §22.6--22.7~\cite{king2008c}}
\end{frame}

\begin{frame}[fragile]{Desafio}
  \vspace{-0.4em}
  \begin{columns}[T]
    \begin{column}{0.50\textwidth}
      \begin{block}{Busca binária no disco}
        \small \texttt{turma.bin} está em ordem crescente de matrícula.
        \texttt{./busca turma.bin mat} encontra o aluno sem ler o arquivo
        inteiro: compare com o registro do \textbf{meio}, descarte a metade
        que não pode conter \texttt{mat} e repita, com \texttt{fseek} +
        \texttt{fread}.
        \vspace{-0.3em}
        \begin{itemize}
          \item \texttt{ini = 0}, \texttt{fim = n - 1}, \texttt{meio = (ini
            + fim) / 2};
          \item imprima quantos \texttt{fread} foram necessários.
        \end{itemize}
      \end{block}
    \end{column}
    \begin{column}{0.46\textwidth}
\begin{lstlisting}[style=shellS, language={}]
$ ./busca turma.bin 202600103
202600103 Carla Menezes Santos (leituras: 2)
$ ./busca turma.bin 202600999
nao encontrado (leituras: 3)
\end{lstlisting}
      \raggedright\small Com um milhão de alunos, quantas leituras, no
      máximo? E lendo do início?
    \end{column}
  \end{columns}
  \fonte{Knuth (1998), §6.2.1, Algoritmo B~\cite{knuth1998art}}
\end{frame}

\begin{frame}{Fechamento}
  \begin{columns}[T]
    \begin{column}{0.5\textwidth}
      \textbf{O que mudou hoje}\\[0.3em]
      \small O arquivo virou um vetor de registros em disco: salta-se
      direto ao aluno $k$ e troca-se uma nota no lugar.
    \end{column}
    \begin{column}{0.5\textwidth}
      \textbf{O preço}\\[0.3em]
      \small Ninguém lê \texttt{turma.bin} num editor, e outra máquina pode
      não entendê-lo. E todo nome ocupa 50 bytes, até \texttt{Ana}.
    \end{column}
  \end{columns}
  \begin{block}{Para levar}
    \begin{itemize}
      \small
      \item \texttt{fwrite}/\texttt{fread} copiam bytes; teste o retorno;
      \item tamanho do registro: \texttt{sizeof}, nunca a soma à mão;
      \item nada de ponteiros dentro do que vai para o disco;
      \item registro $k$: \texttt{fseek(f, k * (long) sizeof r, SEEK\_SET)};
      \item alterar no lugar: \texttt{"r+b"} e um \texttt{fseek} entre ler e
        gravar.
    \end{itemize}
  \end{block}
\end{frame}

\section*{Referências}
\begin{frame}[allowframebreaks]{Referências}
  \nocite{*}
  \bibliographystyle{plain}
  \bibliography{referencias}
\end{frame}

\end{document}
